\chapter{Input Models: State Preparation and Block Encoding}
\label{ch:lec06}

\begin{objectives}
\begin{itemize}[leftmargin=1.4em,nosep]
  \item Explain the state-preparation oracle as the vector input model.
  \item Define block encoding as the matrix input/output model.
  \item Construct block encodings of sparse and structured matrices.
  \item Distinguish Hermitian from general block encodings.
\end{itemize}
\end{objectives}

\todo{Write the introductory paragraph orienting the lecture.}

\section{State-preparation oracles}
\label{sec:lec06-1}
\todo{Draft this section.}

\section{Block encoding: definition and factor}
\label{sec:lec06-2}
\todo{Draft this section.}

\section{Block encoding of $s$-sparse matrices}
\label{sec:lec06-3}
\todo{Draft this section.}

\section{Hermitian block encoding}
\label{sec:lec06-4}
\todo{Draft this section.}

\section{Block encoding beyond the computational basis}
\label{sec:lec06-5}
\todo{Draft this section.}

\exercisesection
\begin{exercise}
\todo{Add exercise.}
\end{exercise}
\begin{exercise}
\todo{Add exercise.}
\end{exercise}
\begin{exercise}
\todo{Add exercise.}
\end{exercise}

\readingsection
\begin{itemize}[leftmargin=1.4em]
  \item Lin Lin, \emph{qasc notes}, Ch.~6.
  \item Lin Lin, \emph{live notes}, Ch.~9.
  \item Dong An, \emph{Quantum Numerical Linear Algebra} talk (input models).
\end{itemize}
